In applications, it typically depends on boundary conditions whether expansion in a cosine or sine series is desirable; see the examples in Chapter~\ref{chap:introduction} and also the comments below on pointwise convergence of Fourier series. The expansion in terms of sines and cosines provided by Corollary~\ref{cor:fourier-complex-hilbert-basis} is typically used for periodic functions.

It is nice to know that Fourier series converge in $L^2$, but this leaves a number of issues. For example:\begin{enumerate}
\item Under what conditions does the Fourier series represent a function in a pointwise sense?
\item Can Fourier series be differentiated term by term? Of course they can be in the sense of distributions, but it is also of interest to know whether the differentiated series converges in $L^2$.
\end{enumerate}

It is known from measure theory that a sequence converging in $L^2$ has a subsequence which converges almost everywhere. For Fourier series it is actually not necessary to take a subsequence; this is a hard theorem which was not proved until 1966. A much more elementary observation is that $\sum_{n=0}^\infty a_n\cos(n\pi x)$ converges uniformly on $[0,1]$ if $\sum_{n=0}^\infty|a_n|$ converges, and using the Cauchy-Schwarz inequality in $\ell^2$, we can see that this is the case if $\sum_{n=1}^\infty|a_n|^2n^\alpha$ converges for any $\alpha>1$ (set $|a_n|=(|a_n|n^{\alpha/2})(n^{-\alpha/2})$; if $\alpha>1$, then the sequence $n^{-\alpha/2}$ is in $\ell^2$).

Now, let $f\in L^2(0,1)$ be such that the derivative of $f$ (in the sense of distributions) is also in $L^2(0,1)$ (we shall study such functions extensively in the section on Sobolev spaces later (Chapter~\ref{chap:sobolev-spaces})). Then $f'$ can be expanded in either a sine or a cosine series:
\begin{equation}\label{eq:fourier-derivative-sine-cosine-expansions}
\begin{aligned}
f'(x)&=\sum_{n=1}^\infty a_n\sin(n\pi x),\\
f'(x)&=\sum_{n=0}^\infty b_n\cos(n\pi x).
\end{aligned}
\end{equation}
By integration, we find
\begin{equation}\label{eq:fourier-antiderivative-cosine-sine-expansions}
\begin{aligned}
f(x)&=\sum_{n=1}^\infty-\dfrac{a_n}{n\pi}\cos(n\pi x)+\alpha,\\
f(x)&=b_0x+\sum_{n=1}^\infty\dfrac{b_n}{n\pi}\sin(n\pi x)+\beta.
\end{aligned}
\end{equation}


The first of these expressions represents a cosine series for $f$, and since \[\sum_{n=1}^\infty|a_n^2|<\infty\] we have \[\sum_{n=1}^\infty|a_n|/n\pi<\infty,\] i.e., the first series in \eqref{eq:fourier-antiderivative-cosine-sine-expansions} converges uniformly. Hence any $f\in L^2(0,1)$ which has a derivative in $L^2(0,1)$ has a uniformly convergent cosine series. (In particular, this implies that any such $f$ is continuous. This is a special case of the Sobolev embedding theorem.) Moreover, in the sense of $L^2$-convergence, the series can be differentiated term by term.

The second series in \eqref{eq:fourier-antiderivative-cosine-sine-expansions}, on the other hand, is a sine series only if $\beta=b_0=0$. It is easy to see that $\beta=f(0)$ and $b_0=\int_0^1 f'(x)\,\mathrm{d}x=f(1)$. Hence any function $f\in L^2(0,1)$ such that $f'\in L^2(0,1)$ and in addition $f(0)=f(1)=0$ has an absolutely convergent sine series. This shows that the convergence behavior of a Fourier series is influenced not only by the smoothness of the function but also by its behavior at the boundary.
